\newpage\handout[true]
{\textsc{Math 2106-B: Foundations of Mathematical Proof}}{Fall 2025}
{Homework 7: More on Groups and Isomorphisms}
{\textit{Prof.} \href{https://ceheitsch.github.io/webpage/}{Dr. Christine Heitsch}}{Due Date: October 30, 2025}
{Math 2106-B: HW7}
{\large
  \noindent
  \textbf{Name:} \HandoutAuthor \smallskip \\

  \noindent
  \textbf{GTID:} 903743506 \smallskip \\

  \noindent
  %%% List anyone with whom you discussed the problems
  \textbf{Collaborators:} None. This material reflects independent preparation. \smallskip \\

  \noindent
  %%% List all resources used OTHER than the textbook, lecture notes, 
  %%% webwork problems, ICA questions, and previous homework assignments.
  \textbf{Outside resources:} This work was informed by MATH 2106 course materials 
  (ICAs: in-class activities and WebWork contents) from \HandoutInstitution\ (the course textbook is 
  Chartrand \parencite{chartrand_mathematical_2018}).   
  Outside references specifically include
  textbooks written by 
  Grimaldi \parencite{grimaldi_discrete_2004} (used in MATH 3012) and 
  Rosen \parencite{rosen_discrete_2019} (used in CS 2050). \smallskip \\

  \noindent
  \textbf{Notice}: The answers contain cross-references throughout this document as clickable hyperlinks in most PDF viewers. However, the Gradescope PDF view might not support clickable hyperlinks.
}
\printbibliography[
  heading=bibintoc,   
  title={References}     % change “References” text if needed
]
% `heading=bibnumbered` if you prefer it numbered chapter (in \tableofcontents)
% `heading=bibintoc`    if you prefer it unnumbered but still in the ToC
\newpage
% \printbibliography[
  heading=bibintoc,   
  title={References}     % change “References” text if needed
]
% `heading=bibnumbered` if you prefer it numbered chapter (in \tableofcontents)
% `heading=bibintoc`    if you prefer it unnumbered but still in the ToC\newpage
% The following table presents the standard logical equivalences from these sources together with 
the author's own (or the student, myself, Jaehoon Song) clarifications and labels where applicable.


\begin{table}[h!]
  \centering
  \renewcommand{\arraystretch}{1.2}
  \setlength{\tabcolsep}{8pt}
  \begin{tabular}{|p{5.0cm}|p{5.0cm}|p{5.2cm}|}
    \hline
    \rowcolor[HTML]{E0E0E0}
    \multicolumn{2}{|l|}{\textbf{Logical Equivalences}} & \textbf{Name} \\
    \hline
    $\,p \AND q \equiv q \AND p$ & $\,p \OR q \equiv q \OR p$ & Commutative Law \\
    \hline
    $\, (p \AND q) \AND r \equiv p \AND (q \AND r)$ & $\, (p \OR q) \OR r \equiv p \OR (q \OR r)$ & Associative Law \\
    \hline
    $\,p \AND (q \OR r) \equiv (p \AND q) \OR (p \AND r)$ & $\,p \OR (q \AND r) \equiv (p \OR q) \AND (p \OR r)$ & Distributive Law \\
    \hline
    $\,\overline{p \AND q} \equiv \sNOT{p} \OR \sNOT{q}$ & $\,\overline{p \OR q} \equiv \sNOT{p} \AND \sNOT{q}$ & DeMorgan's Law \\
    \hline
    \multicolumn{2}{|l|}{$\,\overline{\overline{p}} \equiv p$} & Double Negation law \\
    \hline
    $\,p \AND \boldsymbol{T} \equiv p$ & $\,p \OR \boldsymbol{F} \equiv p$ & Identity Law \\
    \hline
    $\,p \OR \boldsymbol{T} \equiv \boldsymbol{T}$ & $\,p \AND \boldsymbol{F} \equiv \boldsymbol{F}$ & Domination Law \\
    \hline
    $\,p \AND p \equiv p$ & $\,p \OR p \equiv p$ & Idempotent Law \\
    \hline
    \multicolumn{2}{|l|}{$\,p \THEN q \equiv \sNOT{p} \OR q$} & Implication Equivalence \\
    \hline
    \multicolumn{2}{|l|}{$\,\overline{p \THEN q} \equiv p \AND \sNOT{q}$} & Negation of Implication \\
    \hline
    \multicolumn{2}{|l|}{$\,p \THEN q \equiv \sNOT{q} \THEN \sNOT{p}$} & Contrapositive Equivalence \\
    \hline
    \multicolumn{2}{|l|}{$\,p \OR \sNOT{p} \equiv \boldsymbol{T}$} & Law of Tautology \\
    \hline
    \multicolumn{2}{|l|}{$\,p \AND \sNOT{p} \equiv \boldsymbol{F}$} & Law of Contradiction \\
    \hline
  \end{tabular}
\end{table}\newpage
% 
The following definitions are presented from standard mathematical sources (\textbf{Sets}).
Clarifications and labels have been added where applicable.

\begin{define}{Basic Set Concepts}\phantomsection\label{def:basic-set-concepts} % \hyperref[def:basic-set-concepts]{YOUR_CONTEXT}
  \begin{enumerate}
    \item \textbf{Universal Set:} A universal set $\setUNIV$ is a
    set that contains all objects under consideration in a particular context or problem.
    \item \textbf{Empty Set:} The empty set, denoted $\setNULL$,
    is the unique set that contains no elements.
    \item \textbf{Subset:} Let $A$ and $B$ be sets. $A$ is a
    subset of $B$ ($A \subseteq B$) if (and only if)
    for all $x \in A$, $x \in B$.
    \item \textbf{Set Equality:} Let $A$ and $B$ be sets. The sets $A$ and $B$ are
    equal ($A = B$) if $A \subseteq B$ and $B \subseteq A$.
    \item \textbf{Power Set:} Let $A$ be a set. The power set of $A$, denoted $\sPOWER{A}$ or $2^A$,
    is the set of all subsets of $A$. That is, $\mathcal{P}(A) = \{S\in\setUNIV : S \subseteq A\}$.
  \end{enumerate}
\end{define}
The following definitions extend the understanding of set operations
and properties.


\begin{define}{Set Operations}\phantomsection\label{def:set-operations} % \hyperref[def:set-operations]{YOUR_CONTEXT}
  \begin{enumerate}
    \item \textbf{Complementary Set:} Let $A$ be a subset of the universal set $\setUNIV$.
    The complement of $A$, denoted $A^c$ or $\overline{A}$, is the set of all elements
    in $\setUNIV$ that are not in $A$. That is, $A^c = \{x \in \setUNIV : x \notin A\}$.
    \item \textbf{Union:} Let $A$ and $B$ be sets. The union of $A$ and $B$,
    denoted $A \cup B$, is the set of all elements that are in $A$, or in $B$, or in both.
    That is, $A \cup B = \{x : x \in A \lor x \in B\}$.
    \item \textbf{Intersection:} Let $A$ and $B$ be sets. The intersection of $A$ and $B$,
    denoted $A \cap B$, is the set of all elements that are in both $A$ and $B$.
    That is, $A \cap B = \{x : x \in A \land x \in B\}$.
    \item \textbf{Set Difference:} Let $A$, $B$ be subsets of a universal set $\setUNIV$.
    The set difference $A \sMINUS B$ is defined to be the set
    $\{ x \in \setUNIV \LDIV x \in A \text{ and } x \notin B\}$. \vspace*{10px}\\
    \textit{It is sometimes called the relative complement of $B$ in $A$,
    and denoted in the textbook as $A - B$.}
    \item \textbf{Cartesian Product:} Let $A$ and $B$ be sets. The Cartesian product of
    $A$ and $B$, denoted $A \sTIMES B$, is the set consisting of all ordered pairs:
    \[
    A \sTIMES B = \{(a, b) \LDIV a \in A, b \in B\}.
    \]
  \end{enumerate}
\end{define}


\newpage

The following claims are given in the reference as well. Let $A,B\in U$, a universal set.


\begin{claim}{Subset of Union: $A \subseteq \sOR{A}{B}$}\phantomsection\label{claim:subset-union} % \hyperref[claim:subset-union]{YOUR_CONTEXT}
  \begin{subprove}[bluecolor]
      Let $x\in U$, assume $x\in A$. Then the statement \inlinequote{$x\in A$ or $x\in B$} is true.
      Hence, by the definition, $\sOR{A}{B}=\{x\in U\LDIV x\in A$ or $x\in B\}$, $x\in \sOR{A}{B}$.
      Thus, $A \subseteq \sOR{A}{B}$.
  \end{subprove}
\end{claim}

\begin{claim}{Empty Subset: $\setNULL \subseteq A$}\phantomsection\label{claim:empty-subset} % \hyperref[claim:empty-subset]{YOUR_CONTEXT}
  \begin{subprove}[bluecolor]
      Let $x\in U$, assume $x\in \setNULL$. But there are no elements in $\setNULL$, so
      $x\in \setNULL$ is false. Since the premise is false, the implication if $x\in \setNULL$,
      then $x\in A$ is true. Thus, $\setNULL \subseteq A$ by definition of \hyperref[def:basic-set-concepts]{subsets}.
  \end{subprove}
\end{claim}

\begin{claim}{Subset as Equality: $\sOR{A}{B}=A$ if and only if $B\subseteq A$}\phantomsection\label{claim:subset-equality} % \hyperref[claim:subset-equality]{YOUR_CONTEXT}
  \begin{subprove}[bluecolor]
      Let $x\in U$, first, we assume $B\subseteq A$, we have already shown
      $A \subseteq \sOR{A}{B}$ by the claim, \hyperref[claim:subset-union]{Subset of Union}. 
      To show $\sOR{A}{B} \subseteq A$, assume $x\in \sOR{A}{B}$.
      By definition, $x\in A$ or $x\in B$. Considering $x\in B$ and assumption
      $B \subseteq A$, $x\in B$ implies $x\in A$. Thus, $\sOR{A}{B} \subseteq A$.
      Hence, if $B\subseteq A$, then $\sOR{A}{B} = A$. \vspace*{10px}\\
      Now, we assume $\sOR{A}{B}=A$. To show $B\subseteq A$, we now assume $x\in B$. Since $x\in B$, by definition
      of union, $x\in \sOR{A}{B} = A$. Therefore, $B\subseteq A$.
  \end{subprove}
\end{claim}


%%%%%%%%%%%%%%%%%%% 9/8/25

Let $\setUNIV$ be a set. 
\begin{claim}{ICA6}\phantomsection\label{claim:ica6} % \hyperref[claim:ica6]{YOUR_CONTEXT}
  For all $A,B\in\setUNIV$, $$A\sTIMES B =\setNULL\THEN \group{A=\setNULL\OR B=\setNULL}$$
  \begin{subprove}[bluecolor]
    Let $A,B\in\setUNIV$. We prove $A\sTIMES B=\setNULL$. Assume it is not the case that 
    $A=\setNULL$ or $B=\setNULL$. Thus, we have $A=\setNULL$ and $B=\setNULL$. 
    Since $A\ne\setNULL$ by assumption, there exists an element $a\in A$. Likewise, there 
    is some $b\in B$ since $B\ne\setNULL$. Thus, $\group{a,b}\in\sBUILD{(x,y)}{x\in A,y\in B}$ 
    by definition. But then, $A\sTIMES B\ne\setNULL$ by definition. Hence, the claim holds. 
  \end{subprove}
\end{claim}


\begin{define}{Relatively Prime Integers (Coprime)}\phantomsection\label{def:relatively-prime} % \hyperref[def:relatively-prime]{YOUR_CONTEXT}
  Let $x,y \in \mathbb{Z}$. The integers $x$ and $y$ are said to be
  \textbf{relatively prime} if and only if
  $\gcd(x,y) = 1$.
\end{define}
\newpage
% The following definitions and axioms are presented from standard mathematical sources (\textbf{Number Theory}) together with the author's own (or the student, myself, Jaehoon Song) clarifications and labels where applicable.

\begin{define}{Integers}[def:integers][red]
  $\setZ$: the set of all integers
\end{define}

\begin{define}{Rational Numbers}[def:rational-numbers][red]
  $\setQ$: the set of all rational numbers
\end{define}

\begin{define}{Real Numbers}[def:real-numbers][red]
  $\setR$: the set of all real numbers
\end{define}

\begin{define}{Natural Numbers}[def:natural-numbers][red]
  $\setN$ or $\setZ_{>0}$: the set of all natural numbers which are the set of positive integers (not including $0$)
\end{define}

\begin{define}{Positive Real Numbers}[def:positive-real-numbers][red]
  $\setR_{>0}$: the set of all positive real numbers
\end{define}

\begin{define}{Non-negative Real Numbers}[def:non-negative-real-numbers][red]
  $\setR_{\geq 0}$: the set of all non-negative real numbers
\end{define}


\newpage

The following definitions and claims extend the understanding of integers.

Here are the ICA2 Definitions and Axioms as follows.

\begin{define}{Even Integer}[def:even-integer][red]
  An integer $n$ is \underbar{even} if there exists 
  (existential quantification logic $\exists$) an integer $k$ such that
  \[
  n = 2k.
  \]
\end{define}

\begin{define}{Odd Integer}[def:odd-integer][red]
  An integer $n$ is \underbar{odd} if 
  there exists some integer $k$ such that
  \[
  n = 2k + 1.
  \]
\end{define}

\begin{define}{Closure under Addition}[def:closure-addition][red]
  The sum of two integers is again an 
  integer. (Closure under addition)
\end{define}

\begin{define}{Closure under Multiplication}[def:closure-multiplication][red]
  The product of two integers is again 
  an integer. (Closure under multiplication)
\end{define}

\begin{define}{Odd if and only if Not Even}[def:odd-iff-not-even][red]
  An integer is odd if and only if it 
  is not even.
\end{define}

Here are the ICA3 Definitions and Axioms as follows.

\begin{define}{Universal Closure under Addition}[def:universal-closure-addition][red]
  The sum of two integers (Universal quantification for all, for every, for each) is again an integer. (Closure under addition)
\end{define}

\begin{define}{Universal Closure under Multiplication}[def:universal-closure-multiplication][red]
  The product of two integers (Universal quantification for all, for every, for each) is again an integer. (Closure under multiplication)
\end{define}

\begin{define}{Even Integer Square}[def:even-integer-square][red]
  An integer $n$ is even if and only if $n^2$ is even.
\end{define}

\newpage

The following claims are given in the reference as well.

\begin{define}{Sum of Even Integers}[claim:sum-even-integers][red]
  The sum of two even integers is again an even integer.
  \begin{subprove}[red]
    Let $x$ be an integer and $y$ be an integer ($\forall x,y \in \setZ$). 
    Suppose $x$ is even and $y$ is even. Then, by \nameref{def:even-integer}, there exists an integer 
    $k$ and an integer $j$ such that $x=2k$ and $y=2j$. Then, 
    $x+y=\left(2k\right)+\left(2j\right)=2\left(k+j\right)$. By \nameref{def:closure-addition}, since $k$ and 
    $j$ are integers, then so is $\left(k+j\right)$. Hence, by definition, $x+y$ is an 
    even integer.
  \end{subprove}
\end{define}

The following definitions extend the understanding of integers including divisibility.

\begin{define}{Divisibility}[def:divisibility][red]
  An integer $n$ divides an integer $m$, denoted $n \LDIV m$, if there 
  exists an integer $k$ such that $m=n k$. Otherwise, $n \nmid m$.
\end{define}

The following definitions extend the classification of real numbers into rational and irrational numbers.

\begin{define}{Rational Number}[def:rational-number][red]
  A real number $x$ is \underbar{rational} if there exist integers $a$ and $b$ with $b \neq 0$ such that $bx = a$.
  Otherwise, $x$ is \underbar{irrational}.

  \textbf{Note}: Compare the rational number definition with the even/odd axiom \inlinequote{odd iff not even} from \nameref{def:odd-iff-not-even}.
\end{define}

The following claims are given in the reference as well.

\begin{define}{Irrationality of $\sqrt{2}$}[claim:sqrt2-irrational][red]
  $\sqrt{2}$ is irrational.
  \begin{subprove}[red]
    Suppose $\sqrt{2}$ is rational. Then, by \nameref{def:rational-number}, there exists $a,b\in \setZ$ 
    with $b\ne 0$ such that $b\cdot\sqrt{2}=a$. We may further assume that any factors 
    common to $a$ and $b$ have already been removed.\\

    Then, squaring both sides, this yields $2\cdot b^2 = a^2$.
    Since products of integers are again integers by \nameref{def:closure-multiplication}, we have that $a^2,b^2\in \setZ$.
    Thus, by definition, $a^2$ is even. Therefore, by \nameref{def:even-integer-square}, $a$ is even. 
    Thus, there exists $c\in\setZ$ such that $a=2c$. Then
    $$
    2b^2=a^2={2c}^2=4c^2
    $$
    Then, $b^2=2c^2$. Then, as before, $b^2$ is even, so $b$ is even.
    However, $a$ and $b$ both being even contradicts the fact that any factors 
    common to $a$ and $b$ have already been removed.
  \end{subprove}
\end{define}






%%%%%%%%%%%%%%%%%%% New style (refactor model)
\begin{define}{Relatively Prime Integers (Coprime)}[def:relatively-prime][red]
  Let $x,y \in \mathbb{Z}$. The integers $x$ and $y$ are said to be
  \textbf{relatively prime} if and only if
  $\gcd(x,y) = 1$.
\end{define}\newpage
%%%%%%%%%%%%%%%%%%%%%%%%%%%%%%%%%%%%%%%%%%%%%%%%%%%%%%%%%%%%%%%%%%%%%%%
% References Section
%%%%%%%%%%%%%%%%%%%%%%%%%%%%%%%%%%%%%%%%%%%%%%%%%%%%%%%%%%%%%%%%%%%%%%%
\textbf{Remark}: The following definitions (or results) are already proved in class or on previous ICA, 
Wbwk, Hmwk assignments, and contents of the textbook.
\begin{define}{Groups}[def:group][red]
  A \textbf{group} is a nonempty set \(G\) under the (binary) 
  operation \(*: S \times S \to S\), denoted by \((G, *)\), such that
  \begin{itemize}
    \item \textbf{Closure:} For all \(a, b \in G\), \(a * b \in G\).
    \item \textbf{Associativity:} For all \(a, b, c \in G\), \((a * b) * c = a * (b * c)\).
    \item \textbf{Existence of Identity:} There exists an element \(e \in G\) 
    such that for all \(a \in G\), \(e * a = a * e = a\).
    \item \textbf{Existence of Inverse:} For each \(a \in G\), there exists an 
    element \(a^{-1} \in G\) such that \(a * a^{-1} = a^{-1} * a = e\).
  \end{itemize}
\end{define}
Let $(G, *)$ be a group with identity $e \in G$.
\begin{define}{Subgroup}[def:subgroup][red]
  A nonempty subset $H$ of $G$ is called a subgroup of $G$ if (and only if) $(H, *)$is a group. (That is, $H$ itself is a group with respect to the product of $G$.)
\end{define}
\begin{define}{Subgroup Lemma}[lem:subgroup-lemma][red]
  (To be proved on ICA15.) A subset $H \subseteq G$ is a subgroup of $(G, *)$if (and only if) 
  $H$ is nonempty, closed under the product of $G$, and contains inverses, i.e 
  \begin{enumerate}
    \item $H \neq \emptyset$,
    \item for all $a, b \in H$, $a * b \in H$, and
    \item for all $a \in H, a^{-1} \in H$.
  \end{enumerate}
\end{define}
\begin{define}{Exponent Notation}[def:exponent-notation][red]
  For $a \in G$, define $a^0=e$. 
  For $i \in \setZ$ with $i \geq 1$, 
  define $a^i=a * a^{i-1}$ and $a^{-i}=\left(a^{-1}\right)^i$. 
  It follows that $a^n a^m=a^{n+m}$ and $\left(a^n\right)^m=a^{n m}$ for $n, m \in \setZ$.
\end{define}
\begin{define}{Order of an Element}[def:order-of-an-element][red]
  Let $G$ be a finite group. The \textbf{order} of an 
  element $a \in G$, denoted $o(a)$, is the least positive 
  integer $m$ such that $a^m=e$.
\end{define}
\begin{define}{Abelian Group}[def:abelian-group][red]
  A group $(G, *)$ is called an \textbf{abelian group} (or \textbf{commutative group}) if for all $a, b \in G$, 
  \[
    a * b = b * a.
  \]
\end{define}
\begin{define}{Cyclic Subgroup}[def:cyclic-subgroup][red]
  The cyclic subgroup generated by $a \in G$ is the 
  set, $\left(a\right)=\left\{a^i \mid i \in \setZ\right\}$.
\end{define}
\begin{define}{Symmetric Group}[def:symmetric-group][red]
  The symmetric group of degree $n$, denoted $S_n$, is the group of all permutations (bijective maps) of the set $\{1,\dots,n\}$ with composition as the group operation:
  \[
    S_n=\{\sigma:\{1,\dots,n\}\to\{1,\dots,n\}\mid \sigma\ \text{is a bijection}\}
  \]
  Note that $\abs{S_n}=n!$.
\end{define}


%%%%%%%%%%%%%%%%%%%%%%%%%%%%%%%%%%%%%%%%%%%%%%%%%%%%%%%%%%%%%%%%%%%%%%%
% Problems Section
%%%%%%%%%%%%%%%%%%%%%%%%%%%%%%%%%%%%%%%%%%%%%%%%%%%%%%%%%%%%%%%%%%%%%%%
\newpage
Let $G$ be a group with identity $e$.
\begin{enumerate}
  \item Prove or disprove from the definitions:
  \begin{enumerate}
    \item A cyclic group is abelian.
    \begin{define}{Cyclic Group}[def:cyclic-group][red]
      A group $(G, *)$ is called a \textbf{cyclic group} if there exists an element $g \in G$ such that $G = (g) = \{g^i \mid i \in \setZ\}$, 
      where $(g)$ is the cyclic subgroup generated by $g$ as defined in \nameref{def:cyclic-subgroup}. 
      The element $g$ is called a \textbf{generator} of $G$.
    \end{define}
    \begin{subprove}
      Let $(G, *)$ be a cyclic group and let $g \in G$ be a generator of $G$ such that $G = (g) = \{g^i \mid i \in \setZ\}$.
      Let $a, b \in G$ be arbitrary elements.
      Since $G$ is cyclic, there exist integers $m, n \in \setZ$ such that $a = g^m$ and $b = g^n$.
      By the \nameref{def:exponent-notation},
      \[
        a * b = g^m * g^n = g^{m+n} = g^{n+m} = g^n * g^m = b * a.
      \]
      Since $a$ and $b$ were arbitrary elements of $G$, it follows that for all $a, b \in G$, $a * b = b * a$.
      Therefore, by definition of \nameref{def:abelian-group}, $G$ is an abelian group.
    \end{subprove}
    \item An abelian group is cyclic.
    \begin{subdisprove}
      Let $V_4 = \setZ_2 \times \setZ_2$ be the Klein 4-group under coordinate-wise addition, 
      which was shown to be an abelian group in Homework 6.
      The elements of $V_4$ are $\{([0],[0]), ([0],[1]), ([1],[0]), ([1],[1])\}$.
      
      For $V_4$ to be cyclic, there must exist an element $g \in V_4$ such that $(g) = V_4$.
      However, it is observed that:
      \begin{itemize}
        \item The element $([0],[0])$ is the identity and generates only itself.
        \item The element $([0],[1])$ satisfies $([0],[1])^2 = ([0],[0])$, so $( ([0],[1]) ) = \{([0],[0]), ([0],[1])\}$.
        \item The element $([1],[0])$ satisfies $([1],[0])^2 = ([0],[0])$, so $( ([1],[0]) ) = \{([0],[0]), ([1],[0])\}$.
        \item The element $([1],[1])$ satisfies $([1],[1])^2 = ([0],[0])$, so $( ([1],[1]) ) = \{([0],[0]), ([1],[1])\}$.
      \end{itemize}
      Since no element generates all four elements of $V_4$, it is concluded that $V_4$ is not cyclic.
      Thus, $V_4$ is an abelian group that is not cyclic.
      
      Therefore, by counterexample, the statement that an abelian group is cyclic is false.
    \end{subdisprove}
  \end{enumerate}




  \newpage
  \item If $G$ is cyclic, prove that every subgroup of $G$ is cyclic.
  \begin{define}{Well-Ordering Principle}[def:well-ordering][red]
    Every nonempty subset of the positive integers $\setZ^+$ has a least element. 
    That is, if $S \subseteq \setZ^+$ and $S \neq \emptyset$, then there exists $m \in S$ such that $m \leq n$ for all $n \in S$.
  \end{define}
  \begin{define}{Division Algorithm}[def:division-algorithm][red]
    For any integers $n$ and $d \neq 0$, there exist unique integers $q$ (quotient) and $r$ (remainder) such that 
    $n = qd + r$ and $0 \leq r < |d|$. If $d > 0$, then $0 \leq r < d$.
  \end{define}
  \begin{subprove}
    Let $(G, *)$ be a cyclic group with generator $g \in G$ such that $G = (g) = \{g^i \mid i \in \setZ\}$.
    Let $H$ be a subgroup of $G$ for the sake of direct proof.
    
    Two cases are considered.
    
    \textbf{Case 1:} If $H = \{e\}$, where $e$ is the identity element of $G$, then $H$ is cyclic since it is generated by $e$.
    
    \textbf{Case 2:} If $H \neq \{e\}$, then $H$ contains at least one non-identity element.
    Since $H \subseteq G = (g)$ and $H \neq \{e\}$, there exists an integer $k \neq 0$ such that $g^k \in H$.
    
    By the \nameref{def:well-ordering}, let $d$ be the smallest positive integer such that $g^d \in H$.
    It is claimed that $H = (g^d)$.
    
    To show that $H \subseteq (g^d)$, let $h \in H$ be arbitrary.
    Since $h \in H \subseteq G = (g)$, there exists an integer $n \in \setZ$ such that $h = g^n$.
    By the \nameref{def:division-algorithm}, there exist integers $q$ and $r$ with $0 \leq r < d$ such that $n = qd + r$.
    It follows that
    \[
      h = g^n = g^{qd + r} = (g^d)^q * g^r.
    \]
    Since $g^d \in H$ and $H$ is a subgroup, $(g^d)^q \in H$ by closure under the group operation.
    Moreover, since $h \in H$ and $(g^d)^q \in H$, it follows that $g^r = h * ((g^d)^q)^{-1} \in H$ by closure and existence of inverses in $H$.
    
    However, since $d$ is the smallest positive integer such that $g^d \in H$ and $0 \leq r < d$, 
    the only possibility is that $r = 0$.
    Therefore, $h = (g^d)^q \in (g^d)$.
    Since $h$ was arbitrary, it is concluded that $H \subseteq (g^d)$.
    
    To show that $(g^d) \subseteq H$, since $g^d \in H$ and $H$ is a subgroup, 
    it follows by closure under the group operation that $(g^d) = \{(g^d)^i \mid i \in \setZ\} \subseteq H$.
    
    Therefore, $H = (g^d)$, and hence $H$ is cyclic by definition of \nameref{def:cyclic-group}.
    
    In both cases, it is shown that $H$ is cyclic.
    Since $H$ was an arbitrary subgroup of $G$, it follows that every subgroup of $G$ is cyclic.
  \end{subprove}





  \newpage
  
  \item A subgroup is \textit{proper} if it is neither the identity 
  nor the entire group. Prove that if $G$ has no proper subgroups, 
  then $G$ is cyclic. \textit{Do not assume that $G$ is finite}.
  \begin{define}{Proper Subgroup}[def:proper-subgroup][red]
    A subgroup $H$ of $G$ is called a \textbf{proper subgroup} if $H \neq \{e\}$ and $H \neq G$, 
    where $e$ is the identity element of $G$.
  \end{define}
  \begin{subprove}
    Let $(G, *)$ be a group with identity $e$ and suppose that $G$ has no proper subgroups for the sake of direct proof.
    
    Two cases are considered.
    
    \textbf{Case 1:} If $G = \{e\}$, then $G$ is cyclic since it is generated by $e$.
    
    \textbf{Case 2:} If $G \neq \{e\}$, then $G$ contains at least one non-identity element.
    Let $g \in G$ be an arbitrary non-identity element, i.e., $g \neq e$.
    
    Consider the cyclic subgroup $(g) = \{g^i \mid i \in \setZ\}$ generated by $g$, 
    as defined in \nameref{def:cyclic-subgroup}.
    By definition of \nameref{def:subgroup}, $(g)$ is a subgroup of $G$.
    
    Since $g \neq e$, it follows that $(g) \neq \{e\}$.
    Moreover, since $G$ has no proper subgroups, $(g)$ cannot be a proper subgroup of $G$.
    Therefore, it must be that $(g) = G$.
    
    Since $G = (g)$ and $g \in G$, it is concluded by definition of \nameref{def:cyclic-group} that $G$ is cyclic, 
    with $g$ as a generator.
    
    In both cases, it is shown that $G$ is cyclic.
    Since $g$ was an arbitrary non-identity element of $G$ in Case 2, 
    it follows that every non-identity element of $G$ generates $G$.
  \end{subprove}







  
  \item Let $H$ be a nonempty subset of $G$ such that $a b^{-1} \in H$ for all $a, b \in H$. Prove that $H$ is a subgroup of $G$.
  \begin{subprove}
    Let $H$ be a nonempty subset of $G$ such that $a b^{-1} \in H$ for all $a, b \in H$ for the sake of direct proof.
    It is shown that $H$ satisfies the conditions of \nameref{lem:subgroup-lemma}.
    
    \begin{itemize}
      \item \textbf{Nonemptyness:} By hypothesis, $H$ is nonempty.
      
      \item \textbf{Existence of Identity:} Since $H$ is nonempty, there exists an element $a \in H$.
      Taking $a = b$ in the given condition, it follows that $a a^{-1} = e \in H$, where $e$ is the identity element of $G$.
      Therefore, the identity element is contained in $H$.
      
      \item \textbf{Existence of Inverses:} Let $a \in H$ be arbitrary.
      Since $e \in H$ (as shown above) and $a \in H$, taking $a = e$ and $b = a$ in the given condition, 
      it follows that $e a^{-1} = a^{-1} \in H$.
      Therefore, every element in $H$ has its inverse in $H$.
      
      \item \textbf{Closure:} Let $a, b \in H$ be arbitrary.
      Since $a \in H$ and $b^{-1} \in H$ (as shown above), applying the given condition with the first element as $a$ and the second element as $b^{-1}$,
      it follows that $a (b^{-1})^{-1} = a b \in H$.
      Therefore, $H$ is closed under the group operation.
    \end{itemize}
    
    Since $H$ is nonempty, closed under the group operation, and contains inverses, 
    it is concluded by \nameref{lem:subgroup-lemma} that $H$ is a subgroup of $G$.
  \end{subprove}








  \newpage



  \item Let $A$ and $B$ be subgroups of $G$ and let $A B=\{a b \mid a \in A, b \in B\}$.
  \begin{enumerate}
    \item Assume that $G$ is abelian. Prove that $A B$ is a subgroup of $G$.
    \begin{subprove}
      Let $A$ and $B$ be subgroups of $G$ and let $A B = \{a b \mid a \in A, b \in B\}$.
      Assume that $G$ is abelian for the sake of direct proof.
      It is shown that $A B$ satisfies the conditions of \nameref{lem:subgroup-lemma}.
      
      \begin{itemize}
        \item \textbf{Nonemptyness:} Since $A$ and $B$ are subgroups, they both contain the identity element $e$ of $G$.
        Therefore, $e = e e \in A B$, and hence $A B$ is nonempty.
        
        \item \textbf{Closure:} Let $a_1 b_1, a_2 b_2 \in A B$, where $a_1, a_2 \in A$ and $b_1, b_2 \in B$.
        Since $G$ is abelian,
        \[
          (a_1 b_1) (a_2 b_2) = a_1 (b_1 a_2) b_2 = a_1 (a_2 b_1) b_2 = (a_1 a_2) (b_1 b_2).
        \]
        Since $A$ and $B$ are subgroups, $a_1 a_2 \in A$ and $b_1 b_2 \in B$ by closure.
        Therefore, $(a_1 b_1) (a_2 b_2) = (a_1 a_2) (b_1 b_2) \in A B$, and hence $A B$ is closed under the group operation.
        
        \item \textbf{Inverses:} Let $a b \in A B$, where $a \in A$ and $b \in B$.
        Since $A$ and $B$ are subgroups, $a^{-1} \in A$ and $b^{-1} \in B$ by existence of inverses.
        Since $G$ is abelian,
        \[
          (a b)^{-1} = b^{-1} a^{-1} = a^{-1} b^{-1} \in A B.
        \]
        Therefore, every element in $A B$ has its inverse in $A B$.
      \end{itemize}
      
      Since $A B$ is nonempty, closed under the group operation, and contains inverses,
      it is concluded by \nameref{lem:subgroup-lemma} that $A B$ is a subgroup of $G$.
    \end{subprove}
    \item Disprove the statement that $A B$ is a subgroup of $G$.
    \begin{subdisprove}      
      Let $G = S_3$, the symmetric group of degree 3, as defined in \nameref{def:symmetric-group}.
      Let $A = \{e, (12)\}$ and $B = \{e, (13)\}$ be subgroups of $S_3$,
      where $e$ is the identity permutation and $(12)$ and $(13)$ are transpositions.
      
      The set $A B$ is computed as:
      \begin{align*}
        A B &= \{a b \mid a \in A, b \in B\} \\
        &= \{e e, e (13), (12) e, (12) (13)\} \\
        &= \{e, (13), (12), (12)(13)\} \\
        &= \{e, (12), (13), (123)\}.
      \end{align*}
      
      It is claimed that $A B$ is not a subgroup of $S_3$.
      To show this, consider the elements $(12)$ and $(13)$ in $A B$.
      Their product is $(12) (13) = (123)$, which is in $A B$.
      However, consider $(13) (12) = (132)$.
      It is verified that $(132) \notin A B = \{e, (12), (13), (123)\}$.
      
      If $A B$ were a subgroup, then for any $x, y \in A B$, the product $x y$ would be in $A B$.
      However, $(13) (12) = (132) \notin A B$, which contradicts the closure property.
      Thus, $A B$ is not a subgroup of $S_3$.
      
      Therefore, by counterexample, the statement that $A B$ is always a subgroup of $G$ is false.
    \end{subdisprove}
  \end{enumerate}


  
  \newpage


  \item Let $(G, *)$and $(H, \circ)$be groups with identities $e \in G, f \in H$ and $\phi: G \rightarrow H$ an isomorphism.
  \begin{define}{Group Isomorphism}[def:isomorphism][red]
    A function $\phi: G \to H$ between groups $(G, *)$ and $(H, \circ)$ is called an \textbf{isomorphism} if:
    \begin{enumerate}
      \item $\phi$ is a bijection (one-to-one and onto),
      \item $\phi$ preserves the group operation, i.e., $\phi(a * b) = \phi(a) \circ \phi(b)$ for all $a, b \in G$.
    \end{enumerate}
    If such an isomorphism exists, $G$ and $H$ are said to be \textbf{isomorphic}, denoted $G \cong H$.
  \end{define}
  \begin{enumerate}
    \item Prove that $\phi(e)=f$.
    \begin{subprove}
      Let $(G, *)$ and $(H, \circ)$ be groups with identities $e \in G$ and $f \in H$, 
      and let $\phi: G \to H$ be an isomorphism as defined in \nameref{def:isomorphism}.
      
      Since $\phi$ is an isomorphism, it preserves the group operation, i.e., 
      $\phi(a * b) = \phi(a) \circ \phi(b)$ for all $a, b \in G$.
      
      For any $g \in G$, it is observed that $e * g = g$.
      Applying $\phi$ to both sides and using the property that $\phi$ preserves the group operation,
      \[
        \phi(e * g) = \phi(e) \circ \phi(g) = \phi(g).
      \]
      Since $\phi(g) \in H$ and $f$ is the identity in $H$, it follows that
      \[
        \phi(e) \circ \phi(g) = \phi(g) = f \circ \phi(g).
      \]
      By the cancellation property in $H$, it is concluded that $\phi(e) = f$.
    \end{subprove}
    \item Prove that $\phi\left(g^{-1}\right)=(\phi(g))^{-1}$ for all $g \in G$.
    \begin{subprove}
      Let $g \in G$ be arbitrary and let $(G, *)$ and $(H, \circ)$ be groups with identities $e \in G$ and $f \in H$, 
      and let $\phi: G \to H$ be an isomorphism as defined in \nameref{def:isomorphism}.
      
      Since $\phi$ preserves the group operation and by part (a), $\phi(e) = f$,
      it follows that
      \[
        \phi(g * g^{-1}) = \phi(g) \circ \phi(g^{-1}) = \phi(e) = f.
      \]
      Since $\phi(g) \circ \phi(g^{-1}) = f$ and $f$ is the identity in $H$,
      it is concluded that $\phi(g^{-1}) = (\phi(g))^{-1}$.
      
      Since $g$ was arbitrary, it follows that $\phi(g^{-1}) = (\phi(g))^{-1}$ for all $g \in G$.
    \end{subprove}
    \item Let $g \in G$ have order $k \in \mathbb{N}$. Prove that $\phi(g)$ has order $k$ in $H$.
    \begin{subprove}
      Let $g \in G$ have order $k \in \mathbb{N}$ as defined in \nameref{def:order-of-an-element}, 
      meaning that $k$ is the least positive integer such that $g^k = e$.
      Let $(G, *)$ and $(H, \circ)$ be groups with identities $e \in G$ and $f \in H$, 
      and let $\phi: G \to H$ be an isomorphism as defined in \nameref{def:isomorphism}.
      
      First, it is shown that $(\phi(g))^k = f$.
      Since $\phi$ preserves the group operation, it follows by induction that 
      $\phi(g^k) = (\phi(g))^k$ for any positive integer $k$.
      Therefore,
      \[
        (\phi(g))^k = \phi(g^k) = \phi(e) = f,
      \]
      where the last equality follows from part (a).
      
      Next, it is shown that $k$ is the least positive integer with this property.
      Suppose, for contradiction, that there exists a positive integer $m < k$ such that $(\phi(g))^m = f$.
      Then, since $\phi$ preserves the group operation and is bijective,
      \[
        \phi(g^m) = (\phi(g))^m = f = \phi(e).
      \]
      Since $\phi$ is injective (one-to-one), it follows that $g^m = e$.
      However, this contradicts the assumption that $k$ is the least positive integer such that $g^k = e$, 
      since $m < k$ and $g^m = e$.
      
      Therefore, $k$ is the least positive integer such that $(\phi(g))^k = f$.
      By definition of \nameref{def:order-of-an-element}, it is concluded that $\phi(g)$ has order $k$ in $H$.
    \end{subprove}
  \end{enumerate}




  \newpage



  \item Fix $a \in G$. Prove that the function $\phi: G \rightarrow G$ defined by $\phi(g)=a g a^{-1}$ for $g \in G$ is an isomorphism.
  \begin{subprove}
    Let $(G, *)$ be a group with identity $e$ and fix $a \in G$.
    Let $\phi: G \to G$ be defined by $\phi(g) = a g a^{-1}$ for all $g \in G$ for the sake of direct proof.
    It is shown that $\phi$ satisfies the conditions of \nameref{def:isomorphism}.
    
    \begin{itemize}
      \item \textbf{Bijection:}
      \begin{itemize}
        \item \textbf{One-to-one:} Let $g_1, g_2 \in G$ be such that $\phi(g_1) = \phi(g_2)$.
        Then $a g_1 a^{-1} = a g_2 a^{-1}$.
        Multiplying both sides on the left by $a^{-1}$ and on the right by $a$,
        \[
          a^{-1} (a g_1 a^{-1}) a = a^{-1} (a g_2 a^{-1}) a,
        \]
        which simplifies to $g_1 = g_2$.
        Therefore, $\phi$ is injective (one-to-one).
        
        \item \textbf{Onto:} Let $h \in G$ be arbitrary.
        It is claimed that there exists $g \in G$ such that $\phi(g) = h$.
        Taking $g = a^{-1} h a$, it is verified that
        \[
          \phi(g) = \phi(a^{-1} h a) = a (a^{-1} h a) a^{-1} = (a a^{-1}) h (a a^{-1}) = e h e = h.
        \]
        Therefore, $\phi$ is surjective (onto).
      \end{itemize}
      Since $\phi$ is both injective and surjective, it is concluded that $\phi$ is a bijection.
      
      \item \textbf{Operation Preservation:} Let $g_1, g_2 \in G$ be arbitrary.
      Then,
      \[
        \phi(g_1 * g_2) = a (g_1 * g_2) a^{-1} = a g_1 (a^{-1} a) g_2 a^{-1} = (a g_1 a^{-1}) * (a g_2 a^{-1}) = \phi(g_1) * \phi(g_2).
      \]
      Therefore, $\phi$ preserves the group operation.
    \end{itemize}
    
    Since $\phi$ is a bijection and preserves the group operation,
    it is concluded by definition of \nameref{def:isomorphism} that $\phi$ is an isomorphism from $G$ to $G$.
  \end{subprove}



  \newpage
  



  \item Let $(G, *)$be a group. Define a binary operation $\circ$ on $G$ by $a \circ b=b * a$.
  \begin{enumerate}
    \item Prove that $(G, \circ)$is a group.
    \begin{subprove}
      Let $(G, *)$ be a group with identity $e$ and define the binary operation $\circ$ on $G$ by $a \circ b = b * a$ for all $a, b \in G$ for the sake of direct proof.
      It is shown that $(G, \circ)$ satisfies the conditions of \nameref{def:group}.
      
      \begin{itemize}
        \item \textbf{Closure:} Let $a, b \in G$.
        Since $(G, *)$ is a group and $b * a \in G$ by closure,
        it follows that $a \circ b = b * a \in G$.
        Therefore, $(G, \circ)$ is closed under the operation $\circ$.
        
        \item \textbf{Associativity:} Let $a, b, c \in G$.
        Then,
        \[
          (a \circ b) \circ c = (b * a) \circ c = c * (b * a) = (c * b) * a = a \circ (c * b)=a \circ (b \circ c)
        \]
        since $c * (b * a) = (c * b) * a$. It follows that $(a \circ b) \circ c = a \circ (b \circ c)$.
        Therefore, $(G, \circ)$ is associative under the operation $\circ$.
        
        \item \textbf{Existence of Identity:} It is claimed that $e$ is the identity element for $(G, \circ)$.
        For any $a \in G$,
        \[
          e \circ a = a * e = a \quad \text{and} \quad a \circ e = e * a = a,
        \]
        where the equalities follow from the fact that $e$ is the identity in $(G, *)$.
        Therefore, $e$ is the identity element in $(G, \circ)$.
        
        \item \textbf{Existence of Inverse:} Let $a \in G$ be arbitrary.
        Since $(G, *)$ is a group, there exists $a^{-1} \in G$ such that $a * a^{-1} = a^{-1} * a = e$.
        It is claimed that $a^{-1}$ is the inverse of $a$ in $(G, \circ)$.
        Then,
        \[
          a \circ a^{-1} = a^{-1} * a = e \quad \text{and} \quad a^{-1} \circ a = a * a^{-1} = e.
        \]
        Therefore, every element in $(G, \circ)$ has an inverse.
      \end{itemize}
      
      Since $(G, \circ)$ satisfies all four group axioms, it is concluded by definition of \nameref{def:group} that $(G, \circ)$ is a group.
    \end{subprove}
    \item Prove that $(G, *)$and $(G, \circ)$are isomorphic.
    \begin{subprove}
      Let $(G, *)$ be a group with identity $e$ and let $(G, \circ)$ be the group with the operation $a \circ b = b * a$ as shown in part (a).
      Define the function $\phi: (G, *) \to (G, \circ)$ by $\phi(g) = g^{-1}$ for all $g \in G$,
      where $g^{-1}$ denotes the inverse of $g$ in $(G, *)$.
      It is shown that $\phi$ is an isomorphism as defined in \nameref{def:isomorphism}.
      
      \begin{itemize}
        \item \textbf{Bijection:}
        \begin{itemize}
          \item \textbf{One-to-one:} Let $g_1, g_2 \in G$ be such that $\phi(g_1) = \phi(g_2)$.
          Then $g_1^{-1} = g_2^{-1}$.
          Taking inverses on both sides (in $(G, *)$), it follows that $(g_1^{-1})^{-1} = (g_2^{-1})^{-1}$,
          which simplifies to $g_1 = g_2$.
          Therefore, $\phi$ is injective (one-to-one).
          
          \item \textbf{Onto:} Let $h \in G$ be arbitrary.
          Since $h \in G$ and $G$ is a group, $h^{-1} \in G$.
          Taking $g = h^{-1}$, it is verified that
          \[
            \phi(g) = \phi(h^{-1}) = (h^{-1})^{-1} = h.
          \]
          Therefore, $\phi$ is surjective (onto).
        \end{itemize}
        Since $\phi$ is both injective and surjective, it is concluded that $\phi$ is a bijection.
        
        \item \textbf{Operation Preservation:} Let $g_1, g_2 \in G$ be arbitrary.
        Then,
        \[
          \phi(g_1 * g_2) = (g_1 * g_2)^{-1} = g_2^{-1} * g_1^{-1} = g_1^{-1} \circ g_2^{-1} = \phi(g_1) \circ \phi(g_2),
        \]
        where the third equality follows from the definition of $\circ$: $g_1^{-1} \circ g_2^{-1} = g_2^{-1} * g_1^{-1}$.
        Therefore, $\phi$ preserves the group operation.
      \end{itemize}
      
      Since $\phi$ is a bijection and preserves the group operation,
      it is concluded by definition of \nameref{def:isomorphism} that $(G, *)$ and $(G, \circ)$ are isomorphic.
    \end{subprove}
  \end{enumerate}







  \newpage

  \item Complete but do not hand in. Write out the product table and repeat problem 8 from \textbf{Hmwk 6} for each of the following groups.
  \begin{enumerate}
    \item $G=<a, x \mid a^{2}=x^{3}=(a x)^{2}=e>$ where the elements are $\left\{a^{i} x^{j} \mid 0 \leq i \leq 1,0 \leq j \leq 2\right\}$.
    \begin{subanswer}
      
    \end{subanswer}
    \item $G=<a, x \mid a^{2}=x^{3}=(x a)^{2}=e>$ where the elements are $\left\{x^{j} a^{i} \mid 0 \leq i \leq 1,0 \leq j \leq 2\right\}$.
    \begin{subanswer}
      
    \end{subanswer}
    \item $G=<a, x \mid a^{2}=x^{4}=(a x)^{2}=e>$ where the elements are $\left\{a^{i} x^{j} \mid 0 \leq i \leq 1,0 \leq j \leq 3\right\}$.
    \begin{subanswer}
      
    \end{subanswer}
    \item $G=<a, x \mid a^{2}=x^{4}=(x a)^{2}=e>$ where the elements are $\left\{x^{j} a^{i} \mid 0 \leq i \leq 1,0 \leq j \leq 3\right\}$.
    \begin{subanswer}
      
    \end{subanswer}
  \end{enumerate}
  \item Complete but do not hand in.
  \begin{enumerate}
    \item Which groups, if any, from Problem 9 are isomorphic? Practice justifying your answer.
    \begin{subanswer}
      
    \end{subanswer}
    \item What, if any, are the names of and/or notation for the different groups?
    \begin{subanswer}
      
    \end{subanswer}
    \item Write the relation $(a x)^{2}=e$ in a least four different ways.
    \begin{subanswer}
      
    \end{subanswer}
  \end{enumerate}






\end{enumerate}










\newpage

Let $S, T, U$, and $V$ be nonempty sets. There are 7 problems on 2 pages of this ICA. If needed, the definition of a group and the Pigeonhole Principle are stated on the second page.

\begin{enumerate}
  \item Let $f: S \rightarrow T, g: T \rightarrow U$ and $h: U \rightarrow V$ be functions.
\end{enumerate}

Prove any of the following statements not already done before.\\
(a) $h \circ(g \circ f)=(h \circ g) \circ f$.\\
(b) $i_{T} \circ f=f$ and $f \circ i_{S}=f$.\\
(c) If $f$ and $g$ are bijective, then so is $g \circ f$ directly from the definitions.\\
2. Prove that a function $f: S \rightarrow T$ is a bijection if and only if there exists a function $g: T \rightarrow S$ with $g \circ f=i_{S}$ and $f \circ g=i_{T}$.\\
3. The set of all one-to-one mappings from a nonempty set $S$ onto itself is denoted $A(S)$. Prove that $A(S)$ is a group under the operation of function composition. For simplicity, write $g \circ f$ as $g f$ for $f, g \in A(S)$.\\
4. When $|S|=n<\infty$, then $A(S)$ is called the symmetric group of degree $n$, often denoted $S_{n}$. The elements of $S_{n}$ are called permutations of $S$, and $S_{n}$ has $n!$ elements.\\
Consider $i, f, g \in S_{3}$ where $i=\left(\begin{array}{lll}1 & 2 & 3 \\ 1 & 2 & 3\end{array}\right), f=\left(\begin{array}{lll}1 & 2 & 3 \\ 2 & 1 & 3\end{array}\right)$, and $g=\left(\begin{array}{lll}1 & 2 & 3 \\ 2 & 3 & 1\end{array}\right)$. The notation means $f(1)=2, f(2)=1, f(3)=3$ while $g(1)=2, g(2)=3, g(3)=1$.\\
(a) Show that $i, f, g, g^{2}, f g, g f$ are all distinct permutations by considering the symmetries of a triangle (under rotation and reflection) with vertices labeled 1,2 , and 3 .\\
(b) Show that $f^{2}=f f=i$ and $g^{3}=g\left(g^{2}\right)=i$. What are the inverses of the other elements?\\
(c) Complete the product table for $S_{3}$ on the board:

\begin{center}
\begin{tabular}{|c|c|c|c|c|c|c|}
\hline
 & $i$ & $f$ & $g$ & $g^{2}$ & $f g$ & $g f$ \\
\hline
$i$ &  &  &  &  &  &  \\
\hline
$f$ &  & $i$ & $f g$ &  &  &  \\
\hline
$g$ &  & $g f$ &  & $i$ &  &  \\
\hline
$g^{2}$ &  &  &  &  &  &  \\
\hline
$f g$ &  &  &  &  &  &  \\
\hline
$g f$ &  &  &  &  &  &  \\
\hline
\end{tabular}
\end{center}

\begin{enumerate}
  \setcounter{enumi}{4}
  \item Prove that if $|S| \geq 3$, then $A(S)$ is not abelian.
  \item Let $f: S \rightarrow T$ and $g: T \rightarrow U$ be functions.\\
(a) Let $h: T \rightarrow U$ be a function. Prove that if $f$ is surjective and $g \circ f=h \circ f$, then $g=h$.\\
(b) Let $h: S \rightarrow T$ be a function. Prove that if $g$ is injective and $g \circ f=g \circ h$, then $f=h$.
  \item Let $f: S \rightarrow S$ be a function.\\
(a) Suppose $S$ is finite. Prove $f$ is onto if and only if $f$ is one-to-one.\\
(b) Disprove the general statement that $f$ is onto only if $f$ is one-to-one.\\
(c) Disprove the general statement that $f$ is one-to-one only if $f$ is onto.
\end{enumerate}

Definition (Herstein): A nonempty set $G$ is a group under the (binary) operation $*$, denoted $(G, *)$, if (and only if) these four propositions are satisfied:

\begin{enumerate}
  \item Closure: For all $a, b \in G, a * b \in G$.
  \item Associativity: For all $a, b, c \in G,(a * b) * c=a *(b * c)$.
  \item Existence of an identity: There exists $e \in G$ such that $a * e=e * a=a$ for all $a \in G$.
  \item Existence of inverses: For every $a \in G$ there exists $b \in G$ such that $a * b=b * a=e$.
\end{enumerate}

NOTE: We will use this standard definition which explicitly includes the closure property. The textbook's version - where closure is implicit - can cause problems!

The Pigeonhole Principle: if $k+1$ or more objects are placed into $k$ boxes, then there is at least one box containing two or more objects.







\newpage

You may assume that ( $\mathbb{Z},+$ ) is a group with identity 0 , and that multiplication on the integers is closed and associative with identity 1.\\
Also, the following number theory result may be used without proof.\\
If $a, b \in \mathbb{Z}$ with either $a \neq 0$ or $b \neq 0$, then the greatest common divisor of $a$ and $b$ is 1 if and only if there exist $n, m \in \mathbb{Z}$ such that $a n+b m=1$.

\begin{enumerate}
  \item Let $G$ be a group. (When not necessary for clarity the product $a * b$ is written as $a b$ for $a, b \in G$ and the group is denoted as simply $G$ instead of ( $G, *$ ).)\\
(a) Prove that if $a b=a c$, then $b=c$ for all $a, b, c \in G$.\\
(b) Prove that if $b a=c a$, then $b=c$ for all $a, b, c \in G$.
  \item Let $G$ be a group.\\
(a) Prove that $\left(a^{-1}\right)^{-1}=a$ for all $a \in G$.\\
"The inverse of the inverse of $a$ is $a$."\\
(b) Prove that $(a b)^{-1}=b^{-1} a^{-1}$ for all $a, b \in G$.\\
"The inverse of the product of $a$ and $b$ is the product of their inverses in reverse order."
  \item Fix an integer $n \geq 2$ and consider $\mathbb{Z}_{n}$, the integers modulo $n$.
\end{enumerate}

For example, $\mathbb{Z}_{3}=\{[0],[1],[2]\}=\{[-5],[9],[23]\}$, etc. Note, however, $\mathbb{Z}_{3} \neq\{0,1,2\}$ !\\
Define addition on $\mathbb{Z}_{n}$ as $[a]+[b]=[a+b]$ and multiplication as $[a] \cdot[b]=[a b]$ for all $[a],[b] \in \mathbb{Z}_{n}$.\\
(a) Prove that addition on $\mathbb{Z}_{n}$ is well-defined.

Hint: if $a \equiv a^{\prime}(\bmod n)$ and $b \equiv b(\bmod n)$, is $a+b \equiv a^{\prime}+b^{\prime}(\bmod n)$ ?\\
(b) Prove that multiplication on $\mathbb{Z}_{n}$ is also well-defined.\\
4. Prove that $\mathbb{Z}_{n}$ is a group under addition for any $n \geq 2$. Hint: read the directions.\\
5. Write out the product tables for ( $\mathbb{Z}_{n},+$ ) for $n=3,4,5,6$, and identify all their subgroups.\\
6. Write out the product tables for $\mathbb{Z}_{n}^{*}=\mathbb{Z}_{n} \backslash\{[0]\}$ under multiplication for $n=3,4,5,6$. Are these groups? Why (not)?\\
7. Prove that $\mathbb{Z}_{n}^{*}$ is a group under multiplication if and only if $n \geq 2$ is a prime. Hint: directions.




\newpage


Let $(G, *)$ be a group with identity $e$.\\
Exponent notation: For $a \in G$, define $a^{0}=e$. For $i \in \mathbb{Z}$ with $i \geq 1$, define $a^{i}=a * a^{i-1}$ and $a^{-i}=\left(a^{-1}\right)^{i}$. It follows that $a^{n} a^{m}=a^{n+m}$ and $\left(a^{n}\right)^{m}=a^{n m}$ for $n, m \in \mathbb{Z}$.

\begin{enumerate}
  \item Let $H \subseteq G$. Prove that ( $H, *$ ) is a subgroup of ( $G, *$ ) if and only if $H$ is nonempty, closed under the product of $G$, and contains inverses, i.e (1) $H \neq \emptyset$, (2) for all $a, b \in H, a * b \in H$, and (3) for all $a \in H, a^{-1} \in H$.
  \item Let $a \in G$ and consider the set $A=\left\{a^{i} \mid i \in \mathbb{Z}\right\}$. Prove that $A$ is a subgroup of $G$. This proves that (a), the cyclic subgroup generated by $a$, is well-defined.
  \item Consider the group of integers under addition. What is the cyclic subgroup generated by -1 ? What is (2)? Are the odd integers a subgroup of $(\mathbb{Z},+)$ ? Justify your answers.
  \item Does ( $\mathbb{Z}_{4},+$ ) have a proper cyclic subgroup? What about ( $\mathbb{Z}_{5},+$ )? Justify your answers.
  \item Consider the set $\mathbb{Z}_{2} \times \mathbb{Z}_{2}$ under the product of coordinate-wise addition. Prove this is a group. This group, denoted $V_{4}$, is called the Klein 4-group.
  \item A group $G$ is cyclic if there exists $a \in G$ such that $G=(a)$. Is ( $\mathbb{Z},+$ ) a cyclic group? What about ( $\mathbb{Z}_{n},+$ ) for integer $n>2$ ? Justify your answers.
  \item Prove that the Klein 4-group in problem 5 is abelian, but not cyclic.
  \item For each of the four group axioms, give an example of a set with a binary product for which the axiom fails.
  \item Give an example of a nonabelian group. Justify your answer fully, i.e. why is it a group and how is it nonabelian?
\end{enumerate}


\newpage


Let ( $G, *$ ) and ( $H, \circ$ ) be groups with identities $e \in G, f \in H$. Here $\circ$ is a generic product symbol (like $*$ ) and not function composition.\\
A function $\phi: G \rightarrow H$ is an isomorphism from $G$ onto $H$ if (1) $\phi$ is a bijection, and (2) $\phi(a * b)= \phi(a) \circ \phi(b)$ for all $a, b \in G$.\\
When proving an isomorphism, you should write the product symbol explicitly.

\begin{enumerate}
  \item Show that the groups $(\{-1,1\}, \cdot),\left(\mathbb{Z}_{2},+\right),\left(\mathbb{Z}_{3}^{*}, \cdot\right)$ are all pair-wise isomorphic.
  \item Prove or disprove:\\
(a) $\left(\mathbb{Z}_{5}^{*}, \cdot\right)$ is isomorphic to $\left(\mathbb{Z}_{4},+\right)$.\\
(b) $V_{4}$, the Klein 4-group, is isomorphic to to ( $\mathbb{Z}_{4},+$ ).\\
(c) ( $\mathbb{Z}_{6},+$ ) is isomorphic to $S_{3}$.
  \item Let $m>1$ be an integer and $H_{m}=\{m x \mid x \in \mathbb{Z}\}$. Prove that $H_{m}$ is a subgroup of ( $\mathbb{Z},+$ ).
  \item Let $G$ be a group with identity $e$ and $H=\left\{a \in G \mid a^{2}=e\right\}$.\\
(a) Prove that if $G$ is abelian, then $H$ is a subgroup of $G$.\\
(b) Is the statement true if $G$ is nonabelian? Justify your answer.
  \item Let $(G, *)$ be a group and $H$ a nonempty subset of $G$ that is closed under $*$.\\
(a) Prove that if $H$ is finite, then $H$ is a subgroup of $G$.\\
(b) Is $H$ a subgroup of $G$ if $G$ is finite? Justify your answer.\\
(c) What about if $G$ is infinite? Justify your answer.
  \item Up to isomorphism, how many groups of order $n$ are there for $n=1,2,3$ ? What about $n=4$ ? And $n=5$ ? Justify your answers with a product table.
  \item Consider the dihedral group of order 8: $D_{8}=<a, x \mid a^{2}=x^{4}=e, a x a^{-1}=x^{-1}>$. with identity $e$, generators $a$ and $x$, and relations $a^{2}=e, x^{4}=e$, and $a x a^{-1}=x^{-1}$.\\
(a) List all 10 subgroups of $D_{8}$.\\
(b) Which subgroups of $D_{8}$ are isomorphic to the Klein 4-group? Justify your answer.
\end{enumerate}



\newpage
Hint: Write one proof outline and then edit as needed.

\begin{enumerate}
  \item Find a formula for the following summations. Prove your formula is correct using induction.\\
(a) $1+2+3+\ldots+n$ for all integers $n \geq 1$.\\
(b) $1+2+2^{2}+\ldots+2^{n}$ for all nonnegative integers $n$.
  \item For which nonnegative integers $n$ do the following inequalities hold? Prove your answer is correct using mathematical induction.\\
(a) $2 n \leq n^{2}$\\
(b) $n^{2}<2^{n}$\\
(c) $2^{n} \leq n$ !\\
(d) $n!<n^{n}$
  \item Proof or spoof? Justify your answer.
\end{enumerate}

Claim: All horses are the same color. Proof: Let $P(n)$ be the proposition that all horses in a set of $n$ horses are the same color. Clearly, $P(1)$ is true. Now assume that $P(n)$ is true, so that all horses in any set of $n$ horses are the same color. Consider any $n+1$ horses; number these as horses $1,2,3, \ldots, n, n+1$. Now the first $n$ of these horses all must have the same color, and the last $n$ of these must also have the same color. Since the set of the first $n$ horses and the last $n$ horses overlap, all $n+1$ must be the same color. This shows that $P(n+1)$ is true and finishes the proof by induction.\\
4. Prove that $n^{2}-1$ is divisible by 8 whenever $n$ is an odd positive integer.\\
5. (a) Determine which amounts of postage can be formed using just 5 -cent and 6 -cent stamps.\\
(b) Prove your answer is correct using mathematical induction.\\
(c) Prove your answer is correct using the "strong" form of mathematical induction.\\
6. Let $h$ be a real number with $h>-1$. Prove that $1+n h \leq(1+h)^{n}$ for all $n \mathbb{Z}, n>0$. This is called Bernoulli's inequality.\\
7. Fix $a, r \in \mathbb{R}$. A geometric progression is an infinite sequence of the form $a, a r, a r^{2}, \ldots, a r^{n}, \ldots$. Find a formula for the sum of the first $n$ terms of a geometric progression, and prove it is correct using induction.

\newpage

\begin{enumerate}
  \item Let $h$ be a real number with $h>-1$. Prove that $1+n h \leq(1+h)^{n}$ for all $n \in \mathbb{Z}, n>0$. This is called Bernoulli's inequality.
  \item Fix $a, r \in \mathbb{R}$. A geometric progression is an infinite sequence of the form $a, a r, a r^{2}, \ldots, a r^{n}, \ldots$. Find a formula for the sum of the first $n$ terms of a geometric progression, and prove it is correct using induction.
  \item Prove by induction that a set having $n$ elements has exactly $2^{n}$ subsets.
  \item Prove by induction on $k$ that a set of size $n$ has exactly $\binom{n}{k}$ subsets of size $k$ where $\binom{n}{k}= \frac{n!}{k!(n-k)!}$.
  \item Let $m>1$ be an integer and $H_{m}=\{m x \mid x \in \mathbb{Z}\}$. Prove that $H_{m}$ is a subgroup of ( $\mathbb{Z},+$ ).
  \item Let $G$ be a group with identity $e$ and $H=\left\{a \in G \mid a^{2}=e\right\}$.\\
(a) Prove that if $G$ is abelian, then $H$ is a subgroup of $G$.\\
(b) Is the statement true if $G$ is nonabelian? Justify your answer.
  \item Let ( $G, *$ ) be a group and $H$ a nonempty subset of $G$ that is closed under *.\\
(a) Prove that if $H$ is finite, then $H$ is a subgroup of $G$.\\
(b) Is $H$ a subgroup of $G$ if $G$ is finite? Justify your answer.\\
(c) What about if $G$ is infinite? Justify your answer.
\end{enumerate}